\section{Logarithmic and dilogarithmic Borel boundary models}
\label{sec:Borel}

The coefficient theorem gives more information than the Borel radius. It
identifies explicit functions carrying the two leading boundary
singularities. The regularity statements in this section concern limits
from the closed disk. Analytic continuation outside that disk is proved
separately in Section~\ref{sec:continuation}.

Put $z=cu$ and abbreviate the two combinations in \eqref{eq:K} by
\begin{equation}\label{eq:PQ}
 P(u)=du+\frac{1+u}{2}z^2,\qquad Q(u)=(2+u)z.
\end{equation}
Thus $K_\epsilon=P H_\epsilon+QH_\epsilon'$.
All logarithms below mean their Taylor branches at $\xi=0$; the
dilogarithms have the same convention. Define
\begin{align}
 B_{0,u}(\xi)
 &=\frac{2u}{\xi}
 \left\{e^{-z}\log(1-\xi/A)+e^z\log(1+\xi/A)\right\},
 \label{eq:B0-model}\\
 B_{1,u}(\xi)
 &=B_{0,u}(\xi)-\frac{2u}{\xi}
 \left\{(P-Q)e^{-z}\Li_2(\xi/A)
                  +(P+Q)e^z\Li_2(-\xi/A)\right\}.
 \label{eq:B1-model}
\end{align}
The factors $\xi^{-1}$ have removable singularities at zero.

\begin{theorem}[Boundary models with controlled remainders]
\label{thm:Borel-models}
For every $u\in\C$, there are functions $\mathcal H_{0,u}$ and
$\mathcal H_{1,u}$, holomorphic in $|\xi|<A$, such that
\begin{equation}\label{eq:Borel-decompositions}
 \widehat V_u=B_{0,u}+\mathcal H_{0,u}
             =B_{1,u}+\mathcal H_{1,u}.
\end{equation}
The function $\mathcal H_{0,u}$ extends continuously to $|\xi|\le A$.
Both $\mathcal H_{1,u}$ and its complex derivative extend continuously
to that closed disk. These extensions, and the convergent series defining
them, are locally uniform in $u$.

For every $u\ne0$, both $\xi=A$ and $\xi=-A$ are genuine singularities
of $\widehat V_u$, whose Taylor radius is exactly $A$.
\end{theorem}
\begin{proof}
Resolve the hyperbolic functions in Theorem~\ref{thm:main} into their two
exponentials. For all sufficiently large $n$, uniformly for $|u|\le U$,
\begin{align}
 \frac{C_n(u)}{\Gamma(n)}
 ={}&-\frac{2u}{nA^n}
       \left\{e^{-z}+(-1)^n e^z\right\}\notag\\
 &-\frac{2u}{n^2A^n}
       \left\{(P-Q)e^{-z}+(-1)^n(P+Q)e^z\right\}
       +O_U\!\left(\frac{|u|}{A^n n^3}\right).
 \label{eq:two-action-coefficients}
\end{align}
For example, at even $n$ the second brace is
$2(P\cosh z+Q\sinh z)$; at odd $n$ its negative is
$2(P\sinh z+Q\cosh z)$. This verifies the signs without any
singularity-transfer inference.

Now use the convergent identities
\[
 \log(1-w)=-\sum_{n\ge1}\frac{w^n}{n},
 \qquad
 \Li_2(w)=\sum_{n\ge1}\frac{w^n}{n^2},
 \qquad |w|<1.
\]
The first line of \eqref{eq:two-action-coefficients} is precisely the
coefficient sequence of \eqref{eq:B0-model}. The first two lines together
give \eqref{eq:B1-model}. Consequently the coefficients of
$\mathcal H_{0,u}$ at $\xi^{n-1}$ are
$O_U(|u|A^{-n}n^{-2})$, and those of $\mathcal H_{1,u}$ are
$O_U(|u|A^{-n}n^{-3})$. The finitely many omitted indices cause no
difficulty: their coefficients are entire functions of $u$.

On $|\xi|\le A$, the first remainder series is therefore dominated by a
constant times $\sum n^{-2}$. The second remainder and its termwise
derivative are dominated by $\sum n^{-3}$ and $\sum n^{-2}$, respectively.
The Weierstrass criterion proves all the asserted extensions, including
local uniformity in the parameter. Equation~\eqref{eq:two-action-coefficients}
also proves that the original Borel series converges in $|\xi|<A$.

For $0<\varepsilon<1$, \eqref{eq:B0-model} and continuity of its remainder
give
\begin{align}
 \widehat V_u(A(1-\varepsilon))
 &=\frac{2u}{A}e^{-z}\log\varepsilon+O_u(1),\label{eq:radial-positive}\\
 \widehat V_u(-A(1-\varepsilon))
 &=-\frac{2u}{A}e^z\log\varepsilon+O_u(1).\label{eq:radial-negative}
\end{align}
Indeed the opposite logarithm tends to $\log2$, and replacing
$(1-\varepsilon)^{-1}$ by $1$ contributes
$O(\varepsilon|\log\varepsilon|)$. For $u\ne0$, neither displayed
logarithmic coefficient vanishes. Thus neither boundary point admits a
holomorphic extension of the principal germ, and its radius is exactly $A$.
\end{proof}

This argument remains valid when an individual parity amplitude vanishes
at a complex value of $u$. The two boundary amplitudes are exponentials
times $u$, and therefore never vanish for $u\ne0$. At $u=0$ the defining
equation gives $V=0$, so there is no singularity.
The exact positive Borel radius also implies that the ordinary
$q$-series $V(q,u)$ has radius zero for each fixed $u\ne0$:
the root growth of $\Gamma(n)$ converts a positive Borel coefficient
growth rate into an unbounded ordinary coefficient growth rate.

The continuously differentiable remainder makes one further radial term
explicit. With $\sigma\in\{1,-1\}$, set
\[
 S_\sigma(u)=\frac{2u}{\sigma A}e^{-\sigma z}.
\]
There are finite constants $L_\sigma(u)$ such that
\begin{equation}\label{eq:radial-refinement}
\begin{split}
 \widehat V_u\bigl(\sigma A(1-\varepsilon)\bigr)
 ={}&S_\sigma(u)\log\varepsilon+L_\sigma(u)\\
 &+S_\sigma(u)\{1-P(u)+\sigma Q(u)\}
                   \varepsilon\log\varepsilon+O_u(\varepsilon).
\end{split}
\end{equation}
To see this, use
$\Li_2(1-\varepsilon)=\zeta(2)+\varepsilon\log\varepsilon
-\varepsilon+O(\varepsilon^2|\log\varepsilon|)$ in
\eqref{eq:B1-model}, expand $(1-\varepsilon)^{-1}$, and use the bounded
derivative of $\mathcal H_{1,u}$. The dilogarithm with opposite argument
is holomorphic near $-1$. This proves \eqref{eq:radial-refinement} directly.
Neither \eqref{eq:Borel-decompositions} nor
\eqref{eq:radial-refinement} asserts that the remainder extends
holomorphically through the two endpoints.

\section{Resurgent continuation and sectorial sums}
\label{sec:continuation}

The preceding boundary models follow from coefficient estimates alone.
The arithmetic forward kernel supplies the additional analytic information
needed for continuation and summation. We use a standard nonlinear
closure theorem, with its hypotheses stated explicitly.

\subsection{The arithmetic kernel and its continuation paths}

For a formal series $f(q)=\sum_{n\ge1}f_nq^n$, our Borel convention is
$\borel f(\xi)=\sum_{n\ge1}f_n\xi^{n-1}/\Gamma(n)$, as in
\eqref{eq:BorelV}. Let
\[
 \Omega=A\Z.
\]
An $\Omega$-continuable germ is holomorphic at zero and admits analytic
continuation along every path starting there whose subsequent points avoid
$\Omega$. Initial motion inside the regular disk is understood in this
definition. When passing to functions of $x=1/q$, write
\[
 \mathcal R_\Omega
 =\C\oplus\borel^{-1}
     \{\text{$\Omega$-continuable germs}\}.
\]
Thus $\mathcal R_\Omega$ is a space of formal series in $x^{-1}$.
Constants are separated because our Borel transform acts on the part
without constant term.

\begin{proposition}[A normally convergent logarithmic lattice]
\label{prop:forward-Borel}
The shifted Borel transform of \eqref{eq:R} is
\begin{equation}\label{eq:R-Borel-lattice}
 \widehat R(\xi)
 =-\frac14-\frac2\xi
       \sum_{m\ge1}d(m)
           \log\left(1-\frac{\xi^2}{A^2m^2}\right).
\end{equation}
Its principal expression is holomorphic on
$\C\setminus\bigl((-\infty,-A]\cup[A,\infty)\bigr)$.
The series continues normally on compact subsets of the path surface
avoiding $\Omega$. In particular $R(1/x)\in\mathcal R_\Omega$.
\end{proposition}
\begin{proof}
For $|\xi|<A$, expansion of each logarithm and the absolutely convergent
Dirichlet identity $\sum_{m\ge1}d(m)m^{-s}=\zeta(s)^2$ give
\[
 -\frac2\xi\sum_{m\ge1}d(m)
       \log\left(1-\frac{\xi^2}{A^2m^2}\right)
 =\sum_{j\ge1}\frac{4\zeta(2j)^2}{2jA^{2j}}\xi^{2j-1}.
\]
These are exactly the coefficients $r_{2j}/\Gamma(2j)$; the constant
$-1/4$ supplies $r_1$.

Fix a bounded continuation path and a compact neighborhood of its lift.
For all sufficiently large $m$, the argument
$\xi^2/(A^2m^2)$ stays in a fixed disk of radius less than one, so its
logarithm retains the Taylor branch and is $O(m^{-2})$, uniformly on that
compact set. The sum of its majorants converges because
$\sum d(m)m^{-2}=\zeta(2)^2<\infty$. Only finitely many logarithms need
nontrivial continuation; their branch points are the stated lattice points
$\pm Am$. This proves normal continuation. The factor $\xi^{-1}$ is
removable at zero on the principal sheet, since each logarithm starts
quadratically.
\end{proof}

The role of zero in $\Omega$ matters. After a loop around a branch point,
one logarithm can change by $2\pi\ii$, adding a nonzero multiple of
$\xi^{-1}$ to \eqref{eq:R-Borel-lattice}. Thus zero can be singular on
another sheet although the original germ is regular there. Our continuation
statement permits this behavior and does not promise passage through zero
after a nontrivial loop.

\subsection{The nonlinear closure input}

We use the following established consequences of Sauzin's convolution
estimates \cite[Theorems 3 and 5, and their summable counterparts
$3'$ and $5'$]{Sauzin2015}. They are quoted here as analytic input.

\begin{proposition}[Nonlinear resurgence and summability closure]
\label{prop:Sauzin-input}
Suppose $\Omega\subset\C$ is closed, discrete, contains zero, and is stable
under addition. Then $\mathcal R_\Omega$ is an algebra with the following
properties.
\begin{enumerate}
 \item If $f_1,\ldots,f_j\in\mathcal R_\Omega$ have zero constant term
 and $H$ is a convergent power series near the origin in $j$ variables,
 then $H(f_1,\ldots,f_j)\in\mathcal R_\Omega$.
 \item The formal maps $x+\mathcal R_\Omega$ form a group under
 composition. In particular a map $G(x)=x+O(1)$ in this group has a
 compositional inverse there, and $f\circ G\in\mathcal R_\Omega$
 whenever $f\in\mathcal R_\Omega$.
\end{enumerate}
The same operations preserve $1$-summability in an arc of directions.
On sufficiently small common sectors, Borel--Laplace summation commutes
with these operations, including compositional inversion.
\end{proposition}

The topology underlying this statement controls analytic continuation on
compact subsets of the path surface. Its estimates justify infinite
substitution and inverse series; coefficientwise formal existence by itself
would not supply these conclusions. The set $A\Z$ satisfies every
hypothesis of the proposition.

\begin{theorem}[Continuation of the quadratic inverse]
\label{thm:inverse-continuation}
For every fixed $u\in\C$, the series $V(1/x,u)$ belongs to
$\mathcal R_\Omega$. Hence \eqref{eq:BorelV} admits analytic continuation
along every path starting at the origin and thereafter avoiding $A\Z$.
\end{theorem}
\begin{proof}
First take $u\ne0$ and define, as formal series,
\begin{equation}\label{eq:hG}
 h(t)=\frac{\sqrt{1-2u^2R(t)}-1}{u},
 \qquad
 G(x)=x\exp\{-h(1/x)\}.
\end{equation}
The square root has constant term one. It is a convergent function of
$R(t)$ near zero, so Proposition~\ref{prop:Sauzin-input} gives
$h(1/x)\in\mathcal R_\Omega$. It satisfies the exact formal identities
\begin{equation}\label{eq:h-equation}
 h(t)+\frac u2 h(t)^2+uR(t)=0,\qquad
 h(t)=\frac u4t+O(t^2).
\end{equation}
Analytic substitution also puts $e^{-h(1/x)}-1$ in
$\mathcal R_\Omega$, with zero constant term.

We check explicitly the multiplication by $x$ in \eqref{eq:hG}.
If $f(x)=\sum_{n\ge1}f_nx^{-n}$ has Borel transform $\widehat f$, then
\[
 \borel\{xf(x)-f_1\}=\partial_\xi\widehat f(\xi).
\]
Differentiation preserves analytic continuation along the allowed paths.
It follows that
$G(x)=x-u/4+O(x^{-1})$ belongs to $x+\mathcal R_\Omega$.
The closure theorem therefore provides its unique compositional inverse
$T=G^{-1}\in x+\mathcal R_\Omega$.

Set
\begin{equation}\label{eq:inverse-parametrization}
 t(q)=\frac1{T(1/q)},\qquad W(q)=h(t(q)).
\end{equation}
Composition closure gives $W(1/x)=h(1/T(x))\in\mathcal R_\Omega$.
The identity $G(T(x))=x$ becomes $q=t e^{h(t)}$, so
$t=qe^{-W}$. Equation~\eqref{eq:h-equation} now implies
\[
 W+\frac u2W^2+uR(qe^{-W})=0.
\]
Formal uniqueness in Proposition~\ref{prop:lagrange} gives $W=V$.
For $u=0$, equation~\eqref{eq:implicit} gives $V=0$ directly.
\end{proof}

This is a continuation theorem with an explicit allowed set. It does not
assert that every nonzero point of $A\Z$ is singular on every sheet.
The genuine singularities at $\pm A$ on the principal boundary were proved
independently in Theorem~\ref{thm:Borel-models}. No assertion about an
all-orders alien derivative or a complete local singularity type is used.

\subsection{Growth away from the real directions and analytic inversion}

\begin{lemma}[Growth of the forward kernel on nonreal rays]
\label{lem:kernel-growth}
For every $\delta\in(0,\pi/2)$ there is $M_\delta>0$ such that the principal
kernel in \eqref{eq:R-Borel-lattice} satisfies
\begin{equation}\label{eq:kernel-growth}
 |\widehat R(\xi)|
 \le M_\delta\{1+\log^2(2+|\xi|)\}
\end{equation}
when $\delta\le\arg\xi\le\pi-\delta$ or
$-\pi+\delta\le\arg\xi\le-\delta$.
\end{lemma}
\begin{proof}
The elementary divisor estimate
\[
 D(M):=\sum_{m\le M}d(m)
       =\sum_{a\le M}\left\lfloor\frac Ma\right\rfloor
       \le M(1+\log M)
\]
and summation by parts give
$\sum_{m>M}d(m)m^{-2}=O((1+\log M)/M)$.
Put $\rho=|\xi|/A$ and initially suppose $\rho\ge1$.
Split the logarithmic lattice at $m=2\rho$.

For $m>2\rho$, Taylor bounds give
$|\log(1-\xi^2/(A^2m^2))|\le C\rho^2/m^2$, so the total tail before
division by $\xi$ is $O(\rho\log(2\rho))$.
For $m\le2\rho$, the angular restriction implies
\[
 |1-s e^{2\ii\arg\xi}|
       \ge \sin\delta\,(1+s)\qquad(s\ge0).
\]
Together with the upper bound $1+s$ and the bounded argument of the
principal logarithm, this bounds each logarithm by
$C_\delta+2\log(1+\rho/m)$, hence by $O_\delta(\log(2+\rho))$.
The estimate for $D(2\rho)$ bounds the head by
$O_\delta(\rho\log^2(2+\rho))$.
Division by $|\xi|=A\rho$ proves \eqref{eq:kernel-growth}.
The remaining bounded part of each closed angular sector is controlled
by holomorphy and the removable value at zero.
\end{proof}

\begin{theorem}[Sectorial sums satisfy the inverse equation]
\label{thm:inverse-summability}
Fix $u\in\C$. The series $V(q,u)$ is $1$-summable in every direction
$\theta\notin\pi\Z$. Let
\[
 R_\theta(t)=\int_0^{e^{\ii\theta}\infty}
                 e^{-\xi/t}\widehat R(\xi)\dd\xi
\]
on its Laplace sector, and write $V_\theta=\SU_\theta V$ for the
corresponding sum of the inverse series. For sufficiently small $q$
in any smaller common sector,
\begin{equation}\label{eq:summed-inverse}
 V_\theta(q,u)+\frac u2V_\theta(q,u)^2
       +uR_\theta\bigl(qe^{-V_\theta(q,u)}\bigr)=0.
\end{equation}
This is the unique small solution with $V_\theta(q,u)=O(q)$.
\end{theorem}
\begin{proof}
Choose an open arc about $\theta$ whose closure avoids the real
directions. Proposition~\ref{prop:forward-Borel} and
Lemma~\ref{lem:kernel-growth} give the analytic continuation and exponential
growth condition for $1$-summability of $R$ throughout this arc.
The logarithmic bound is stronger than the required bound of exponential
type.

Apply the summable version of Proposition~\ref{prop:Sauzin-input} to the
same operations as in \eqref{eq:hG}--\eqref{eq:inverse-parametrization}.
Multiplication by $x$ causes no difficulty: its Borel action is
differentiation, and Cauchy's estimate on a slightly smaller angular
domain preserves an exponential growth bound. Thus $h$, $G-x$, $T-x$,
and $V$ are $1$-summable. Summation commutes with their analytic
substitutions and inverse relations on smaller sectors. In particular,
putting
\[
 h_\theta(t)=\frac{\sqrt{1-2u^2R_\theta(t)}-1}{u}
\]
with the branch tending to zero, the summed parametrization has
$V_\theta=h_\theta(t)$ and $t=qe^{-V_\theta}$. This proves
\eqref{eq:summed-inverse}. At $u=0$, all statements follow from $V=0$.

For uniqueness, the summed forward function and its derivative satisfy
$R_\theta(t)=-t/4+O(t^2)$ and
$R_\theta'(t)=-1/4+O(t)$ on every smaller sector.
The map
\[
 w\longmapsto-\frac u2w^2-uR_\theta(qe^{-w})
\]
therefore has derivative $O_u(|q|)$ on a disk $|w|\le M|q|$ and maps
that disk into itself for a suitable fixed $M$ and sufficiently small
$q$. The contraction argument proves uniqueness there. Any solution
that is $O(q)$ enters such a disk after decreasing the sector radius.
\end{proof}

The preceding sum is defined by Borel--Laplace summation of the divergent
formal tail. It is not an ordinary sum of its exponential sectors.
The real-direction completions of the forward Lambert function are
specified in \cite{incoming}; their upper and lower choices can subsequently
be transported through \eqref{eq:summed-inverse}. The continuation theorem
does not require identifying those choices with an arithmetic average.

